% BORRADOR (15-oct, reencuadre). Reemplaza a Tesis/Plantilla/FrontMatter/F04-Abstract.tex.
% Pendiente de revisión del autor. Borrar estas líneas al copiarlo.
% Presentación abreviada y precisa del contenido de un documento, sin agregar interpretación o crítica; no debe exceder una página.
\section*{Resumen}

\pendiente{Cifras provisionales: juez de relevancia local (Llama 3.1 8B) y validación humana del juez pendiente. Se actualizan con la corrida final.}

Muchas plataformas de empleo buscan comparando las palabras de la consulta con las de la oferta, y fallan cuando ambas describen el mismo cargo con términos distintos. Este trabajo evalúa en qué medida un pipeline de generación aumentada por recuperación (RAG) mejora la recuperación de ofertas frente a la búsqueda por palabras clave, qué componentes explican esa mejora y con qué costo en latencia. El caso de estudio es JobAI, una plataforma de búsqueda de empleo remoto preexistente. Siguiendo Design Science Research, se diseñó un protocolo de evaluación reproducible y se aplicó en dos iteraciones: una de corrección del artefacto y otra que incorporó filtros estructurados extraídos de la consulta, búsqueda híbrida con BM25, reordenamiento y un modelo de \textit{embeddings} multilingüe. La evaluación usó un corpus congelado de 1.068 ofertas reales y 100 consultas en cinco categorías, en inglés y en español; comparó la recuperación con BM25 y con la búsqueda por subcadena de la plataforma, midió el aporte de cada componente agregándolos uno a uno y obtuvo los juicios de relevancia de un modelo de lenguaje. Se midieron además la fidelidad de las respuestas, la latencia por etapa y el compromiso entre \textit{recall} y latencia de pgvector y Qdrant con hasta 500.000 vectores.

La configuración completa obtuvo un nDCG@10 de 0,80, frente a 0,56 de BM25 y 0,06 de la búsqueda por subcadena. Partiendo de la recuperación densa (0,68), los filtros aportaron +0,05 y el reordenamiento +0,08, ambos significativos con la prueba de aleatorización pareada y la corrección de Holm; la búsqueda híbrida no mejoró el resultado. La ventaja de la recuperación densa sobre BM25 fue mayor en las consultas de desajuste de vocabulario. En las consultas con restricciones, la proporción de los cinco primeros resultados que las cumplía pasó de 0,41 a 0,97, y en español la configuración completa alcanzó 0,80, frente a 0,35 de BM25. El reordenamiento dominó el costo: 811 de los 835~ms que tomó construir el contexto. A la escala del corpus real, ambos motores devolvieron resultados de la misma calidad; con 500.000 vectores y un filtro del 1\,\%, pgvector con su índice HNSW no superó un \textit{recall} de 0,90 y solo con la búsqueda exacta forzada alcanzó 1,0, en 9~ms, mientras que Qdrant obtuvo 0,98 en 1,7~ms. La evaluación reveló además 12 fallas del artefacto, 10 de ellas silenciosas: el sistema seguía respondiendo con fluidez mientras ignoraba las restricciones de la consulta, contaba mal las ofertas remotas o no usaba el índice vectorial. Evaluar un sistema RAG exige, por tanto, verificar sus componentes y no solo la respuesta final.

\textbf{Palabras clave}: generación aumentada por recuperación, búsqueda de empleo, recuperación de información, búsqueda híbrida, reordenamiento, bases de datos vectoriales, evaluación, Design Science Research.

\newpage
\section*{Abstract}

\pendiente{Provisional figures: local relevance judge (Llama 3.1 8B), human validation of the judge pending. They will be updated with the final run.}

Many job platforms search by matching the words of the query against those of the posting, and fail when both describe the same position in different terms. This work evaluates to what extent a retrieval-augmented generation (RAG) pipeline improves job retrieval over keyword search, which components explain that improvement, and at what latency cost. The case study is JobAI, a pre-existing remote job search platform. Following Design Science Research, a reproducible evaluation protocol was designed and applied in two iterations: one that corrected the artifact and another that added structured filters extracted from the query, hybrid search with BM25, reranking, and a multilingual embedding model. The evaluation used a frozen corpus of 1,068 real job postings and 100 queries in five categories, in English and Spanish; it compared retrieval against BM25 and against the platform's substring search, measured the contribution of each component by adding them one at a time, and obtained relevance judgments from a language model. It also measured answer faithfulness, per-stage latency, and the recall--latency trade-off of pgvector and Qdrant with up to 500,000 vectors.

The full configuration reached an nDCG@10 of 0.80, against 0.56 for BM25 and 0.06 for substring search. Starting from dense retrieval (0.68), filters added +0.05 and reranking +0.08, both significant under a paired randomization test with Holm correction; hybrid search did not improve the result. The advantage of dense retrieval over BM25 was largest for vocabulary-mismatch queries. For queries with constraints, the share of the top five results that met them rose from 0.41 to 0.97, and in Spanish the full configuration reached 0.80, against 0.35 for BM25. Reranking dominated the cost: 811 of the 835~ms needed to build the context. At the scale of the real corpus, both engines returned results of the same quality; with 500,000 vectors and a 1\,\% filter, pgvector with its HNSW index did not exceed a recall of 0.90 and reached 1.0 only with forced exact search, in 9~ms, while Qdrant obtained 0.98 in 1.7~ms. The evaluation also revealed 12 failures in the artifact, 10 of them silent: the system kept answering fluently while ignoring the query's constraints, miscounting remote jobs or not using the vector index. Evaluating a RAG system therefore requires verifying its components and not only the final answer.

\textbf{Keywords}: retrieval-augmented generation, job search, information retrieval, hybrid search, reranking, vector databases, evaluation, Design Science Research.
